% GENERATED FILE. Edit canonical lessons, then run npm run generate:books.
% canonical-source-hash: fnv1a64:b755a5f81cb84179
% canonical-lessons: ZH-C144-zhongyao, ZH-C144-rongyi, ZH-C144-jiandan, ZH-C144-tebie, ZH-C144-juede

\chapter{Easy and Important}
\label{ch:zh-easy-and-important}
\hlchaptermodality{144}

\begin{chapteropening}
\textbf{By the end of this chapter:} \emph{I can say a thing is important, easy, simple or special, and what I think of it.}

Complete the last lesson of chapter 144: I can say a thing is important, easy, simple or special, and what I think of it.
\end{chapteropening}

\section[\texorpdfstring{\zh{重要} (zhòngyào)}{重要 (zhòngyào)}]{\zh{重要} (zhòngyào) — important}

\label{lesson:ZH-C144-zhongyao}



\begin{quote}
Before the new one: say the Mandarin for punctual, then the Mandarin for a long while.
\end{quote}

\subsection*{You'll want to know: \zh{重要}}
\textbf{\zh{重要}} — \emph{zhòngyào} — \textquotedblleft{}important\textquotedblright{}.

\begin{quote}
\textbf{\zh{每天喝水很重要}} \emph{měitiān hē shuǐ hěn zhòngyào} — Drinking water every day is important.
\end{quote}

\begin{cousinweb}
\textbf{\zh{重要}} (\emph{zhòngyào}) is built from \zh{重} 'heavy' and \zh{要}, glossed 'to bind, to agree', both inherited from Middle Chinese, \emph{drjowngX} and \emph{'jiewH} (Guangyun, in Baxter's transcription). How the pair came to mean important is not settled.
\end{cousinweb}

\begin{grammarlens}[title={What you've built}]
One more word for describing things.
\end{grammarlens}

\subsection*{Your turn}
\begin{itemize}
\raggedright
  \item \emph{Say it:} \emph{zhòngyào}
  \item \emph{Say it:} \emph{zhòngyào}, once more
  \item \emph{Say it:} \emph{měitiān hē shuǐ hěn zhòngyào} — Drinking water every day is important.
  \item \emph{Recall:} say the Mandarin for punctual, then the Mandarin for a long while, then say \emph{zhòngyào} again
\end{itemize}

\begin{tcolorbox}[breakable,colback=teal!4,colframe=teal!35!black,title={Before you move on}]
What does \zh{重要} mean? (\textquotedblleft{}Important\textquotedblright{}.) Say it once more.
\end{tcolorbox}
\section[róngyì]{róngyì — easy}

\label{lesson:ZH-C144-rongyi}



\begin{quote}
Before the new one: say the Mandarin for a long while, then the Mandarin for important.
\end{quote}

\subsection*{You'll want to know: róngyì}
\emph{róngyì} — \textquotedblleft{}easy\textquotedblright{}.

\begin{quote}
\emph{hànzì bù róngyì} — Chinese characters aren't easy.
\end{quote}

Its characters come later; for now it is a word for the ear.

\begin{cousinweb}
\emph{róngyì} is built from \emph{róng} 'to hold' and \emph{yì} 'easy', both inherited from Middle Chinese, \emph{yowng} and \emph{yeH} (Guangyun, in Baxter's transcription); the dictionary glosses the second 'hard and easy' and 'simple'.
\end{cousinweb}

\begin{grammarlens}[title={What you've built}]
One more word for describing things.
\end{grammarlens}

\subsection*{Your turn}
\begin{itemize}
\raggedright
  \item \emph{Say it:} \emph{róngyì}
  \item \emph{Say it:} \emph{róngyì}, once more
  \item \emph{Say it:} \emph{hànzì bù róngyì} — Chinese characters aren't easy.
  \item \emph{Recall:} say the Mandarin for a long while, then the Mandarin for important, then say \emph{róngyì} again
\end{itemize}

\begin{tcolorbox}[breakable,colback=teal!4,colframe=teal!35!black,title={Before you move on}]
What does róngyì mean? (\textquotedblleft{}Easy\textquotedblright{}.) Say it once more.
\end{tcolorbox}
\section[jiǎndān]{jiǎndān — simple}

\label{lesson:ZH-C144-jiandan}



\begin{quote}
Before the new one: say the Mandarin for important, then the Mandarin for easy.
\end{quote}

\subsection*{You'll want to know: jiǎndān}
\emph{jiǎndān} — \textquotedblleft{}simple\textquotedblright{}.

\begin{quote}
\emph{zhège hěn jiǎndān} — This one's simple.
\end{quote}

Its characters come later; for now it is a word for the ear.

\begin{cousinweb}
\emph{jiǎndān} is built from \emph{jiǎn} 'brief' and \emph{dān} 'single', both inherited from Middle Chinese, \emph{keanX} and \emph{tan} (Guangyun, in Baxter's transcription); the dictionary glosses the first as a slip of bamboo for writing on, and 'brief'.
\end{cousinweb}

\begin{grammarlens}[title={What you've built}]
One more word for describing things.
\end{grammarlens}

\subsection*{Your turn}
\begin{itemize}
\raggedright
  \item \emph{Say it:} \emph{jiǎndān}
  \item \emph{Say it:} \emph{jiǎndān}, once more
  \item \emph{Say it:} \emph{zhège hěn jiǎndān} — This one's simple.
  \item \emph{Recall:} say the Mandarin for important, then the Mandarin for easy, then say \emph{jiǎndān} again
\end{itemize}

\begin{tcolorbox}[breakable,colback=teal!4,colframe=teal!35!black,title={Before you move on}]
What does jiǎndān mean? (\textquotedblleft{}Simple\textquotedblright{}.) Say it once more.
\end{tcolorbox}
\section[tèbié]{tèbié — special; especially}

\label{lesson:ZH-C144-tebie}



\begin{quote}
Before the new one: say the Mandarin for easy, then the Mandarin for simple.
\end{quote}

\subsection*{You'll want to know: tèbié}
\emph{tèbié} — \textquotedblleft{}special; especially\textquotedblright{}.

\begin{quote}
\emph{zhège dìfang hěn tèbié} — This place is special.
\end{quote}

Its characters come later; for now it is a word for the ear.

\begin{cousinweb}
\emph{tèbié} is built from \emph{tè} 'alone' and \emph{bié} 'different', both inherited from Middle Chinese, \emph{dok} and \emph{bjet} (Guangyun, in Baxter's transcription).
\end{cousinweb}

\begin{grammarlens}[title={What you've built}]
One more word for describing things.
\end{grammarlens}

\subsection*{Your turn}
\begin{itemize}
\raggedright
  \item \emph{Say it:} \emph{tèbié}
  \item \emph{Say it:} \emph{tèbié}, once more
  \item \emph{Say it:} \emph{zhège dìfang hěn tèbié} — This place is special.
  \item \emph{Recall:} say the Mandarin for easy, then the Mandarin for simple, then say \emph{tèbié} again
\end{itemize}

\begin{tcolorbox}[breakable,colback=teal!4,colframe=teal!35!black,title={Before you move on}]
What does tèbié mean? (\textquotedblleft{}Special; especially\textquotedblright{}.) Say it once more.
\end{tcolorbox}
\section[juéde]{juéde — to find (that); to feel; to think (that)}

\label{lesson:ZH-C144-juede}



\begin{quote}
Before the new one: say the Mandarin for simple, then the Mandarin for special.
\end{quote}

\subsection*{You'll want to know: juéde}
\emph{juéde} — \textquotedblleft{}to find (that); to feel; to think (that)\textquotedblright{}.

\begin{quote}
\emph{wǒ juéde hànyǔ hěn yǒuyìsi} — I find Chinese interesting.
\end{quote}

Its characters come later; for now it is a word for the ear.

\begin{cousinweb}
\emph{juéde} is built from \emph{jué}, glossed 'to understand' and 'to wake', and \emph{de}, from a word glossed 'to get', both inherited from Middle Chinese, \emph{kaewk} and \emph{tok} (Guangyun, in Baxter's transcription).
\end{cousinweb}

\begin{grammarlens}[title={What you've built}]
One more word for describing things.
\end{grammarlens}

\subsection*{Your turn}
\begin{itemize}
\raggedright
  \item \emph{Say it:} \emph{juéde}
  \item \emph{Say it:} \emph{juéde}, once more
  \item \emph{Say it:} \emph{wǒ juéde hànyǔ hěn yǒuyìsi} — I find Chinese interesting.
  \item \emph{Recall:} say the Mandarin for simple, then the Mandarin for special, then say \emph{juéde} again
\end{itemize}

\begin{tcolorbox}[breakable,colback=teal!4,colframe=teal!35!black,title={Before you move on}]
What does juéde mean? (\textquotedblleft{}To find (that); to feel; to think (that)\textquotedblright{}.) Say it once more.
\end{tcolorbox}
